\DocumentMetadata{
  pdfversion=2.0,
  pdfstandard=ua-2,
  lang=en-US,
  tagging=on,
  tagging-setup={math/setup=mathml-SE,tabsorder=structure}
}
\documentclass[11pt,letterpaper]{article}
\usepackage[margin=0.68in,includefoot]{geometry}
\usepackage{newcomputermodern}
\usepackage{amsmath,amscd,mathtools}
\usepackage{tikz,adjustbox}
\providecommand{\square}{\mdlgwhtsquare}
\usepackage[hidelinks]{hyperref}
\hypersetup{
  pdftitle={Algebra First-Year Exam, August 2024, Part 1},
  pdfauthor={University of Florida Department of Mathematics},
  pdfsubject={Graduate mathematics examination},
  pdfdisplaydoctitle=true,
  bookmarksopen=true
}
\setcounter{secnumdepth}{0}
\setlength{\parindent}{0pt}
\setlength{\parskip}{0.28em}
\setlength{\emergencystretch}{3em}
\setlength{\leftmargini}{2.15em}
\setlength{\leftmarginii}{2.65em}
\setlength{\leftmarginiii}{3em}
\pagestyle{empty}
\raggedbottom
\allowdisplaybreaks
\NewTaggingSocketPlug{block/list/label}{examproblem}{%
  \tagstructbegin{tag=Lbl}\tagstructend%
  \tagstructbegin{tag=\LBody}%
  \tagstructbegin{tag=\ProblemHeadingTag,title-o={\ProblemHeadingText}}%
  \tagmcbegin{tag=\ProblemHeadingTag,actualtext={\ProblemHeadingText}}%
  #2%
  \tagmcend%
  \tagstructend%
}
\newenvironment{examproblems}[2]{%
  \def\ProblemHeadingTag{#2}%
  \AssignTaggingSocketPlug{block/list/label}{examproblem}%
  \begin{enumerate}%
  \setcounter{enumi}{#1}%
  \renewcommand{\labelenumi}{\textbf{\theenumi.}}%
  \setlength{\topsep}{0.35em}%
  \setlength{\partopsep}{0pt}%
  \setlength{\itemsep}{0.48em}%
  \setlength{\parsep}{0.18em}%
}{\end{enumerate}}
\newenvironment{examsubparts}{%
  \AssignTaggingSocketPlug{block/list/label}{default}%
  \begin{enumerate}%
  \setlength{\topsep}{0.18em}%
  \setlength{\partopsep}{0pt}%
  \setlength{\itemsep}{0.16em}%
  \setlength{\parsep}{0.1em}%
}{\end{enumerate}}
\newenvironment{examinstructions}{%
  \par\begingroup\itshape
}{%
  \par\endgroup\vspace{0.18em}
}
\makeatletter
\renewcommand\subsection{\@startsection{subsection}{2}{0pt}%
  {-1.25ex plus -.25ex minus -.1ex}%
  {0.55ex plus .1ex}%
  {\normalfont\large\bfseries}}
\renewcommand{\@maketitle}{%
  \newpage\null\vskip -1em%
  {\footnotesize\bfseries
    \href{https://www.ufl.edu/}{University of Florida}, \href{https://math.ufl.edu/}{Department of Mathematics}\par}%
  \vskip -0.5em%
  \tagstructbegin{tag=H1,title={Algebra First-Year Exam, August 2024, Part 1}}%
  \tagpdfparaOff%
  \tagmcbegin{tag=H1}%
    {\LARGE\bfseries\href{https://gma.math.ufl.edu/exams/algebra-first-year/}{Algebra First-Year Exam}, August 2024, Part 1\par}%
  \tagmcend%
  \tagpdfparaOn%
  \tagstructend%
  \vskip 0.25em%
  \tagmcbegin{artifact}%
    \hrule height 1pt%
  \tagmcend%
  \vskip 1em%
}
\makeatother
\title{Algebra First-Year Exam, August 2024, Part 1}
\author{}
\date{}
\begin{document}
\maketitle
\thispagestyle{empty}
\begin{examinstructions}
Answer four problems. (If you turn in more, the first four will be graded.) Within reason, you may quote theorems as long as you state them clearly.
\end{examinstructions}
\begin{examproblems}{0}{H2}
\def\ProblemHeadingText{Problem 1.}
\item
Give an example of each, with a few words of explanation.
\begin{examsubparts}
\item[\textnormal{(a)}]
A finite abelian simple group.
\item[\textnormal{(b)}]
A nonabelian simple group.
\item[\textnormal{(c)}]
A finite nilpotent group which is not a~\(p\)-group.
\item[\textnormal{(d)}]
A finite group which is solvable but not nilpotent.
\end{examsubparts}
\def\ProblemHeadingText{Problem 2.}
\item
Consider the symmetric group \(S_6\), and the corresponding alternating group \(A_6\).
\begin{examsubparts}
\item[\textnormal{(a)}]
How many elements of order 4 are there in \(S_6\)? Justify your answer.
\item[\textnormal{(b)}]
How many elements of order 4 are there in \(A_6\)? Justify your answer.
\end{examsubparts}
\def\ProblemHeadingText{Problem 3.}
\item
Let~\(G\) be a finite~\(p\)-group and suppose that~\(G\) acts on a finite set~\(S\). Denote by \(S^G\) the set of all elements of~\(S\) fixed by every element of~\(G\). Show that \(|S| \equiv |S^G| \pmod p\), where \(|A|\) denotes the cardinality of the set~\(A\).
\def\ProblemHeadingText{Problem 4.}
\item
Let~\(G\) be a group of order \(32 \cdot 11 \cdot 41\).
\begin{examsubparts}
\item[\textnormal{(a)}]
Show that every 41-Sylow subgroup of~\(G\) is normal.
\item[\textnormal{(b)}]
Show that if~\(H\) is a group of order \(32 \cdot 11\) then every 11-Sylow subgroup of~\(H\) is normal.
\item[\textnormal{(c)}]
Show that~\(G\) has a normal cyclic subgroup of order \(451 = 11 \cdot 41\).
\end{examsubparts}
\def\ProblemHeadingText{Problem 5.}
\item
Let~\(G\) be a group and let \(g \in G\).
\begin{examsubparts}
\item[\textnormal{(a)}]
Define what is meant by the order of~\(g\).
\item[\textnormal{(b)}]
Suppose that~\(g\) has infinite order, and suppose that~\(N\) is a finite normal subgroup of~\(G\). Prove, from your definition, that the element \(gN \in G/N\) has infinite order.
\item[\textnormal{(c)}]
Suppose that~\(g\) has finite order, suppose that~\(N\) is a finite normal subgroup of~\(G\), and suppose that \(|N|\) is relatively prime to the order of~\(g\). Prove, from your definition, that the element \(gN \in G/N\) has the same order as the element \(g \in G\).
\end{examsubparts}
\def\ProblemHeadingText{Problem 6.}
\item
Show that for any \(n \geq 5\), the alternating group \(A_n\) is the only nontrivial proper normal subgroup of the symmetric group \(S_n\).
\end{examproblems}
\end{document}
